% RSE: Introduction (reviews S9, P1, P7, L gap statement). Every number is a \metric{} macro.
% Citations: keys in refs/references.bib and refs/verification.csv only, each claim limited to what
% refs/notes/<key>.md supports and logged in refs/claims.csv (rows RSE01 onward). No results are previewed here except the shelf
% gradient, which motivates the problem and is defined in Section 3.5.

Closed seas and saline lakes are shrinking \citep{wurtsbaugh2017,yao2023}. Closed basins worldwide lost water storage between 2002 and 2016, and the Caspian basin alone accounts for over half of the loss in Central Eurasia \citep{wang2018}. Warming raises evaporation from lakes, and where precipitation and inflow do not compensate, lake levels and extents fall \citep{woolway2020}. Water levels in inland seas and lakes are projected to fall through this century \citep{prange2020}. On a shallow, flat floor each metre of fall uncovers a wide band of seabed, and where that band lies decides which harbours, channels and habitats lose the sea, including those of the endemic Caspian seal and of sturgeon \citep{court2025}. Exposed lake beds also become sources of dust, as on the dried bed of the Aral Sea \citep{indoitu2015} and at Great Salt Lake \citep{skiles2018}, and Lake Urmia has followed the Aral Sea's course \citep{aghakouchak2015}. Projections of the emerging seabed combine a water-level pathway with the shape of the basin floor, which recent Caspian assessments took from a global bathymetry grid \citep{court2025,duan2025}.

The Caspian has long served as a natural laboratory of coastal response to level change, and the nearshore slope sets how far its shoreline moves \citep{kaplin1995}. On a nearly flat shelf, a small vertical error in a grid therefore moves its predicted shoreline a long way. In the North Caspian, one metre of vertical error moves the shoreline by about \metricdp{rev.w.r.km_per_m_min}{1}--\metricdp{rev.w.r.km_per_m_max}{1}~km (Section~\ref{sec:sep}). The same problem is known on land, where the vertical error of the global elevation models used in exposure studies changes estimates of coastal flood exposure several-fold \citep{kulp2019,hooijer2021}. Global bathymetry grids are interpolated between sparse soundings, and many cells hold no direct measurement \citep{weatherall2015}. From 2010 until at least the 2014 release, GEBCO's Caspian grid was a compilation of Russian hydrographic charts, shifted by a constant, with no chart datum named \citep{gebco2014_doc}. In the current GEBCO release, most of the North Caspian is coded either as depth predicted from satellite gravity or as a pre-generated grid of mixed sources, and its vertical reference there is not documented \citep{gebco2026}. ETOPO1 took its Caspian depths from the contours of a single navigation map \citep{amante2009}.

The Caspian Sea rose rapidly from 1978 to 1995 and has fallen since 1996 \citep{lahijani2023}. The reversal is explained mainly by increased evaporation over the sea \citep{chen2017}. In the altimetry record its annual mean level dropped from \metric{rev.w.levels.annual_1995_bhs_m}~m in 1995 to \metric{rev.w.levels.annual_2025_bhs_m}~m in 2025 in the Baltic Height System, a fall of \metric{rev.w.levels.annual_fall_m}~m (Section~\ref{sec:levels}). Over that time the shoreline retreated across much of the North Caspian shelf (Fig.~\ref{fig:rse_thennow}). One study using \metric{lit.hoseini2025.n_gcms} climate models gives median changes from 2021 to 2100 of \metric{lit.hoseini2025.medium_median_m}~m under a medium-emission scenario and \metric{lit.hoseini2025.high_median_m}~m under a high-emission one \citep{hoseini2025}. Another puts the fall over 2020--2100 at about \metric{lit.samant2023.ssp245_m}~m under SSP2-4.5 and \metric{lit.samant2023.ssp585_m}~m under SSP5-8.5, each including about \metric{lit.samant2023.trend_2100_m}~m from the observed trend carried forward \citep{samant2023}. In a water-balance model forced by the highest-emission scenario with high water extraction, the northern Caspian dries out by mid-century \citep{koriche2021}. Most of the seabed that emerges will lie in the North Caspian, which holds less than \metric{lit.akbari2020.nc_volume_share_pct}\% of the sea's volume \citep{akbari2020} and is only \metric{rev.w.v.depth2026_m_min}--\metric{rev.w.v.depth2026_m_max}~m deep on average at the 2026 level.

% Figure thennow: generated by figures/erl_fig_then_now.py (the opening figure of the ERL Letter);
% placed where the Introduction first describes the fall.
\begin{figure}[!tbp]
  \centering
  \includegraphics[width=\linewidth]{../figures/out/erl_fig_then_now.pdf}
  \caption{\textbf{Three decades of retreat of the North Caspian Sea.}
  (a) Landsat true-colour composite of June--September \metric{erl.landsat.tm_1996.year}, near the high stand ($n = \metric{erl.landsat.tm_1996.n_scenes}$ scenes).
  (b) The same for \metric{erl.landsat.oli_2025.year} ($n = \metric{erl.landsat.oli_2025.n_scenes}$ scenes), with the observed shoreline at four basin levels, from the Landsat waterlines binned by level (key: level, and years of the acquisitions); dashed white line, the study domain.
  (c) Daily basin level from satellite altimetry (DAHITI reference); grey bars, the two composite summers; coloured bars, the acquisitions of each level in (b).}
  \label{fig:rse_thennow}
\end{figure}


One assessment, based on GEBCO\_2022, found that a decline of \metric{lit.court2025.decline_lo_m}--\metric{lit.court2025.decline_hi_m}~m would reduce existing marine protected area coverage by up to \metric{lit.court2025.mpa_loss_max_pct}\% \citep{court2025}. Another took the sea's topography from GEBCO and concluded that a fall to \metric{lit.duan2025.level_2100_m}~m would remove the northern Caspian \citep{duan2025}. Neither study reports a test of its grid against the seabed that has already emerged.

Satellite records have long been used to map the floors of water bodies whose level changes. Waterlines mapped at known water levels give intertidal topography \citep{mason1995,mason1997,heygster2010,sagar2017,bishoptaylor2019}, and the per-pixel frequency of wetting now gives intertidal elevation through time at continental scale \citep{bishoptaylor2026}. On lakes, a waterline at a known level is likewise an isobath of the bed \citep{feng2011,schroder2022}. Lake hypsometry has been derived from flood frequency or water occurrence paired with level series \citep{cretaux2016,busker2019,fassoniandrade2020,armon2020}. Reservoir and global lake bathymetry have been rebuilt from Landsat water extents and satellite altimetry \citep{getirana2018,li2020,huang2025}. These reconstructions are reliable only within the range of levels already observed: below it they extrapolate, with larger errors \citep{li2020}, or give no estimate \citep{huang2025}. Where global elevation models have been tested on a shrinking lake, the test was of aggregate storage or area \citep{huang2026,bocharov2026}. For the Caspian, ETOPO1 was checked against the emerged band only qualitatively \citep{akbari2020}. A whole-sea comparison found satellite areas within \metric{lit.bocharov2026.max_area_diff_pct}\% of those from historical area--level estimates and a chart-digitised elevation model, none of them a global grid \citep{bocharov2026}. In a whole-sea total, errors of opposite sign on different coasts can cancel. Caspian shoreline change has been mapped from satellite images on the southern coast \citep{toorani2021}, on the northeastern coast \citep{kamza2023,krylenko2026} and over the whole sea since 2016 \citep{dedova2026a}. None of these studies assigns heights to the mapped shorelines or tests a bathymetry grid against them. One compared SRTM contours with topographic maps last updated in 1980 and found SRTM unsuitable for tracking the change \citep{kamza2023}. To our knowledge, the marine bathymetry compilations behind emergence projections for inland seas have not been tested, coast by coast, against where the shoreline actually went.

The Caspian's fall makes that test possible. A Landsat waterline dated and levelled by satellite altimetry is a depth contour of seabed that has since emerged. We flood five global grids from the open sea at each observed level and compare their shorelines with these waterlines. We separate grid error from observation error with a noise floor taken from independent halves of the imagery, a control made from the observations themselves, and a budget of systematic level errors. We apply the same audit to Great Salt Lake, Utah, where a lidar surface of the emerged lake bed shows whether the audit can pass a good surface.

We tested three hypotheses, each specified before any test metric was computed (Section~\ref{sec:prereg}):
\begin{description}
  \item[H1.] Global bathymetry grids misplace the North Caspian shorelines observed since the start of the altimetry record, and their errors are largest in the deltaic and vegetated sectors.
  \item[H2.] An elevation model built from levelled waterlines is more accurate than the grids against independent ICESat-2 heights.
  \item[H3.] Correcting the grids for their hindcast error changes near-term emergence, \metric{h3.scenario.0p5_m}--\metric{h3.scenario.2_m}~m below the 2026 level, by a decision-relevant amount.
\end{description}
We also ask whether an offset fitted to flooded area corrects the shoreline, how much of each grid's error exceeds the errors of the observations, and how much seabed the next metres of fall will expose from the observed 2026 shoreline.
